\section{A regularized minimum and its Hessian}\label{sec:minimization}

We argue by contradiction. This section constructs a strictly negative
minimum with enough regularity to support the later operator identities.
It then extracts an entropy-Hessian estimate at that minimum.

\subsection{The auxiliary ports and the functional}

Suppose Theorem~\ref{thm:epni} fails. The endpoint transmissions give
equality, so $0<\eta<1$. By Lemma~\ref{lem:energy-compactness}, normalized
number cutoffs of the inputs converge with bounded energies and convergent
entropies, as do their outputs. Thus there is a fixed strictly violating
pair $(\rho_1^{\mathrm{w}},\rho_2^{\mathrm{w}})$ with finite number support.
Write $a_i=g^{-1}(S(\rho_i^{\mathrm{w}})/n)$. At least one $a_i$ is positive.

Besides the two variable ports $1,2$, introduce a fixed thermal port
$D$, in state $\rho_D=\tau_{r_D}$ with $r_D>0$. Its mixing weight and those
of the variable inputs are
\[
 0<\lambda_D<1,\qquad
 \lambda_1=(1-\lambda_D)\eta,\qquad
 \lambda_2=(1-\lambda_D)(1-\eta).
\]
The retained port, denoted by $0$, has annihilators
$d_{0,j}=\sum_{s=1,2,D}\sqrt{\lambda_s}\,d_{s,j}$ before reduction.
For positive reference parameters $r_1,r_2$, set
\[
 r_0=\lambda_1r_1+\lambda_2r_2+\lambda_Dr_D,
 \qquad h_s=h(r_s)\quad(s=0,1,2,D).
\]
Let $\gamma$ be the retained state obtained from
$\rho_1\otimes\rho_2\otimes\rho_D$. We also replace a small fraction of
that output by its reference thermal:
\[
 \rho_0=(1-\kappa)\gamma+\kappa\tau_{r_0},\qquad
 0<\kappa<1,\qquad w_i=(1-\kappa)\lambda_i\quad(i=1,2).
\]
Unless otherwise indicated, sums indexed by $i$ run over $1,2$.
Since $h(r)=g'(r)$, the quotient $(S(\rho)-ng(r))/h(r)$ expresses the
entropy deviation in the linearized scale of total photon number at the
thermal reference.
For $\varepsilon,\zeta>0$ we minimize, over input states of finite fourth
moment, the functional
\begin{equation}\label{eq:functional}
 \begin{split}
 J_\zeta(\rho_1,\rho_2)
 ={}&\frac{S(\rho_0)-ng(r_0)}{h_0}
 -\sum_i w_i\frac{S(\rho_i)-ng(r_i)}{h_i}\\
 &+\varepsilon\sum_i\frac{w_i}{h_i}
      D(\rho_i\Vert\tau_{r_i})
 +\zeta\sum_i\Tr\rho_iW_i^4.
 \end{split}
\end{equation}

The normalization of the relative-entropy penalty also matches the closing
stationarity calculation. The thermal entropy formula and
\eqref{eq:thermal-relative-entropy} give
\[
 \frac{D(\rho\Vert\tau_r)}{h(r)}
 =-\frac{S(\rho)-ng(r)}{h(r)}+\bigl(\Tr\rho N-nr\bigr).
\]
Thus the input terms pair an entropy coefficient $1+\varepsilon$ with an
energy coefficient $\varepsilon$. Their centered-energy coefficients
differ by one in stationarity, matching the output centered-energy
relation while retaining the strict defect and mean terms.

The parameters are chosen in the following order. First choose $r_i>0$
equal or sufficiently close to $a_i$. At
$\lambda_D=\kappa=\varepsilon=\zeta=0$, the value at the fixed witness
approaches
\[
 \frac{S(\rho_C^{\mathrm{w}})
       -ng(\eta a_1+(1-\eta)a_2)}
      {h(\eta a_1+(1-\eta)a_2)}<0.
\]
The input differences vanish in this limit. If $a_i=0$, this follows
from $g(r)/h(r)\to0$ as $r\downarrow0$; the output denominator has a
strictly positive, finite limit. Fix $r_D>0$ and then choose
$\lambda_D>0$ sufficiently small to retain negativity. Indeed, the
additional mixing can be performed after the original beam splitter;
as $\lambda_D\downarrow0$, its unitary converges strongly on each total
number sector. The resulting states converge in trace norm with bounded
energy, and their entropies converge by Lemma~\ref{lem:energy-compactness}.

Next choose $\varepsilon>0$ small enough to retain negativity and to have
\begin{equation}\label{eq:slack}
 (1+\varepsilon)\sum_i\lambda_ir_i\le r_0,
 \qquad
 (1+\varepsilon)\sum_i\lambda_i(r_i+1)\le r_0+1.
\end{equation}
At $\varepsilon=0$ the gaps are respectively $\lambda_Dr_D$ and
$\lambda_D(r_D+1)$, both positive. Finally take $\kappa>0$ sufficiently
small. The additional bound on $\kappa$ used in Section~\ref{sec:metric}
depends only on the parameters already fixed, so it is compatible with
this choice. The witness has finite fourth moments and finite thermal
relative entropies. Consequently there are constants $c>0$ and
$\zeta_0>0$ such that
\begin{equation}\label{eq:negative-witness}
 J_\zeta(\rho_1^{\mathrm{w}},\rho_2^{\mathrm{w}})\le -c
 \qquad(0<\zeta\le\zeta_0).
\end{equation}
All parameters except $\zeta$ remain fixed from now on.

\subsection{A Gibbs regularity lemma}

The fourth-moment penalty yields more regularity than its definition
alone suggests. We record the needed form argument explicitly.

\begin{lemma}\label{lem:gibbs-regularity}
Let $W=1+N$ on the $n$-mode Fock space. Suppose $k>0$ and $B$ is a
bounded positive operator. The positive form
\[
 K=kW^4+W^{1/2}BW^{1/2},\qquad \mathcal D(K^{1/2})=\mathcal D(W^2),
\]
defines a positive self-adjoint operator with compact inverse. Its
normalized Gibbs state $\omega=e^{-K}/\Tr e^{-K}$ is faithful and
satisfies $\Tr\omega W^7<\infty$. It uniquely minimizes
$\Tr\sigma K-S(\sigma)$ among states with finite fourth moment.
\end{lemma}

\begin{proof}
Put $C=kI+W^{-3/2}BW^{-3/2}$. The form is
$\|C^{1/2}W^2\phi\|^2$ on $\mathcal D(W^2)$, and is closed because
$kI\le C\le\|C\|I$. The operator
\[
 U=W^{-2}C^{-1}W^{-2}
\]
is positive, compact, and injective. Choose a complete orthonormal
eigenbasis $\psi_j$ with $U\psi_j=u_j^{-1}\psi_j$ and $u_j>0$.
Every $\psi_j$ belongs to $\mathcal D(W^2)$ and satisfies the form
eigenvalue equation for $K$ with eigenvalue $u_j$.
For completeness, the vectors
$\sqrt{u_j}\,C^{-1/2}W^{-2}\psi_j$ are orthonormal and complete:
orthonormality follows from $U$, and completeness from the dense range
of $C^{-1/2}W^{-2}$. Parseval's identity therefore gives
\[
 \|C^{1/2}W^2\phi\|^2
   =\sum_j u_j|\langle\psi_j,\phi\rangle|^2
 \qquad(\phi\in\mathcal D(W^2)).
\]
Closedness and finite spectral sums show that the domain of this
diagonal form is exactly $\mathcal D(W^2)$. This constructs $K=U^{-1}$.

The bound $U\le k^{-1}W^{-4}$ gives
\[
 \#\{j:u_j\le a\}
 \le\operatorname{rank}\mathbf1_{\{W^4\le a/k\}},
\]
by comparing dimensions with the complementary number subspace.
The right side grows polynomially in $a$. Also
$\|W^2\psi_j\|^2\le u_j/k$. Test the form eigenvalue equation against
each number vector and multiply that coordinate equation by $W^{-1/2}$.
It gives
\[
 kW^{7/2}\psi_j
   =u_jW^{-1/2}\psi_j-BW^{1/2}\psi_j.
\]
The right side is square summable, so $\psi_j\in\mathcal D(W^{7/2})$,
with the explicit bound
\[
 \|W^{7/2}\psi_j\|
 \le k^{-1}\bigl(u_j+\|B\|\sqrt{u_j/k}\bigr).
\]
Polynomial eigenvalue counting and the exponential weights $e^{-u_j}$
now prove both $0<\Tr e^{-K}<\infty$ and
$\Tr(e^{-K}W^7)<\infty$. Finally, for a state $\sigma$ with finite
fourth moment, all terms in
\[
 \Tr\sigma K-S(\sigma)
   =D(\sigma\Vert\omega)-\log\Tr e^{-K}
\]
are finite. Nonnegativity and the equality case of relative entropy
give the asserted unique minimum.
\end{proof}

\subsection{Attainment and the exact logarithm identity}

\begin{proposition}\label{prop:regularized-minimum}
For every $0<\zeta\le\zeta_0$, the functional $J_\zeta$ attains a
minimum of value at most $-c$. The input energies at these minima are
bounded uniformly in $\zeta$. At each minimum all four states
$\rho_s$, $s=0,1,2,D$, are faithful, have finite seventh moment, and
satisfy $0\le H_s:=-\log\rho_s\le C_sW_s^4$ as forms. In fact
$H_0\le CW_0$. The inputs obey the exact identities
\begin{equation}\label{eq:log-identity}
 \frac{(1+\varepsilon)w_i}{h_i}H_i
 =\frac{1-\kappa}{h_0}T_iH_0
   +\varepsilon w_iN_i+\zeta W_i^4+c_iI,
 \qquad i=1,2,
\end{equation}
where $T_i$ is the output-observable slice with the other two inputs
fixed, and $c_i$ is a real scalar. The moment and logarithm bounds need
not be uniform as $\zeta\downarrow0$.
\end{proposition}

\begin{proof}
Write $E_i=\Tr\rho_iN_i$. Expanding the thermal relative entropy
expresses $J_\zeta$, up to constants depending only on the fixed
parameters, as
\begin{equation}\label{eq:coercive-functional}
 \frac{S(\rho_0)}{h_0}
 -\sum_i\frac{(1+\varepsilon)w_i}{h_i}S(\rho_i)
 +\varepsilon\sum_i w_iE_i
 +\zeta\sum_i\Tr\rho_iW_i^4.
\end{equation}
The first term is nonnegative. The bound
$S(\rho_i)\le ng(E_i/n)=O(\log(1+E_i))$ shows that the positive linear
energy terms dominate every negative entropy term. Thus bounded upper
sublevels have bounded input energies, uniformly in
$0<\zeta\le\zeta_0$. For fixed $\zeta>0$ their fourth moments are
bounded as well. Take a minimizing sequence in the sublevel supplied
by \eqref{eq:negative-witness}. Lemma~\ref{lem:energy-compactness}
gives a trace-norm convergent subsequence of each input. The product
and channel maps are trace-norm continuous, and output energies are
bounded by the combined input energy plus the fixed replacement
energy. Entropy is continuous along this subsequence. The energy and
fourth-moment expectations are lower semicontinuous, by bounded
spectral truncation. Equation~\eqref{eq:coercive-functional} therefore
gives an admissible minimum, and the same sublevel estimate gives its
uniform energy bound.

Fix such a minimum. The replacement gives $\rho_0\ge\kappa\tau_{r_0}$.
The resolvent representation of the logarithm, together with inverse
order for positive operators, implies
\[
 0\le H_0\le-\log(\kappa\tau_{r_0})\le C W_0
\]
as positive forms. To define $T_i$, pull an output observable back
through the passive unitary and take its partial expectation in the
states of the other two ports. For bounded $Z$, this is a bounded
unital positive slice satisfying
\begin{equation}\label{eq:slice-expectation}
 \Tr\rho_iT_iZ=\Tr\gamma Z.
\end{equation}
For $H_0$, use increasing bounded spectral truncations to define the
positive form $T_iH_0$. Since the retained number is bounded by total
number before reduction, slicing $H_0\le CW_0$ yields
\[
 0\le T_iH_0\le C_iW_i.
\]
In particular, its expectation is the output cross entropy for every
varied input of finite energy.

Vary only input $i$ and replace the output entropy in $J_\zeta$ by
the cross entropy $\Tr\widetilde\rho_0H_0$. Nonnegativity of relative
entropy makes the resulting functional $\overline J_i$ a majorant,
with equality at the chosen minimum. Thus
\[
 \overline J_i(\widetilde\rho_i)
 \ge J_\zeta(\widetilde\rho_i)
 \ge J_\zeta(\rho_i)=\overline J_i(\rho_i).
\]
After division by its entropy coefficient
$(1+\varepsilon)w_i/h_i$, this majorant is
$\Tr\widetilde\rho_iK_i-S(\widetilde\rho_i)$ plus a constant, where
\begin{equation}\label{eq:gibbs-hamiltonian}
 K_i=\frac{h_i}{(1+\varepsilon)w_i}
 \left(\frac{1-\kappa}{h_0}T_iH_0
       +\varepsilon w_iN_i+\zeta W_i^4\right).
\end{equation}
This is a form $kW_i^4+R_i$ with $k>0$ and
$0\le R_i\le C_iW_i$. Hence
$R_i=W_i^{1/2}B_iW_i^{1/2}$ for a bounded positive $B_i$.
Lemma~\ref{lem:gibbs-regularity} identifies the input uniquely as
\[
 \rho_i=\frac{e^{-K_i}}{\Tr e^{-K_i}}.
\]
It is faithful, has finite seventh moment, and satisfies
$H_i=K_i+\log\Tr e^{-K_i}$, proving \eqref{eq:log-identity} and the
claimed input logarithm bounds.

The fixed thermal port has every number moment. Total-number
preservation and
$(x_1+x_2+x_D)^7\le3^6(x_1^7+x_2^7+x_D^7)$ for nonnegative numbers
give finite seventh moment of the raw output; the thermal replacement
preserves this property. This proves all remaining assertions.
\end{proof}

\subsection{The component entropy-Hessian estimate}

Work at one of the minima in Proposition~\ref{prop:regularized-minimum}.
For each label $s$, fix an eigenbasis
$\rho_s=\operatorname{diag}(p_{s,l})$, with all $p_{s,l}>0$. These orthonormal eigenbases generally differ from the number bases used to define the energy, and they depend on the minimizing states. With
$\ell(a,b)=(a-b)/(\log a-\log b)$ and $\ell(a,a)=a$, define
\begin{equation}\label{eq:logmean-metric}
 \|Z\|_s^2
   =\sum_{l,r}h_s\ell(p_{s,l},p_{s,r})|Z_{lr}|^2,
 \qquad
 v_s(l,r)=\frac{\log(p_{s,l}/p_{s,r})}{h_s}.
\end{equation}
Up to the factor $h_s$, this is the Bogoliubov inner product on
observables. Its dual form on state perturbations is the
Bogoliubov--Kubo--Mori metric, obtained from the Hessian of relative
entropy \cite{PetzToth1993,LesniewskiRuskai1999}. We establish the needed
differentiation formulas on finite eigenblocks and then use density;
no bounded inverse for the full state is assumed.

These weighted matrix spaces are Hilbert spaces, with inner product
conjugate-linear in the first variable. Bounded observables belong to
them because the logarithmic mean is at most the arithmetic mean.
The centered subspace is the orthogonal complement of $I$; the
orthogonal projection subtracts
$\langle Z\rangle_sI$, where $\langle Z\rangle_s=\Tr\rho_sZ$.
Set
\begin{equation}\label{eq:slice-map}
 L_iZ=\frac{T_iZ-\langle T_iZ\rangle_i I}{\sqrt{\lambda_i}}.
\end{equation}
Initially the centered output tests are finite eigenblock matrices
with their expectation times $I$ subtracted. They form a dense
subspace. Such a test can have an infinite identity tail, but has
only finitely many frequencies $v_0$. Multiplication by any function
of $v_0$ preserves this class and its centering: every diagonal
entry, including that tail, has frequency zero.

\begin{proposition}\label{prop:hessian}
For $i=1,2$, the slice map on the centered output space satisfies
\begin{equation}\label{eq:hessian}
 \|L_iZ\|_i^2
 \le\frac{1+\varepsilon}{1-\kappa}\,\|Z\|_0^2.
\end{equation}
It consequently extends continuously to the completed centered
spaces. For bounded observables this extension agrees with
\eqref{eq:slice-map}.
\end{proposition}

\begin{proof}
For a single varied input of finite fourth moment, denote the
resulting modified output by $\widetilde\rho_0$. The logarithm
identity \eqref{eq:log-identity} cancels all linear terms in the
objective difference, giving exactly
\begin{equation}\label{eq:relative-entropy-difference}
 J_\zeta(\widetilde\rho_i)-J_\zeta(\rho_i)
 =\frac{(1+\varepsilon)w_i}{h_i}
        D(\widetilde\rho_i\Vert\rho_i)
  -\frac1{h_0}D(\widetilde\rho_0\Vert\rho_0)\ge0.
\end{equation}
All cross entropies are finite by the fourth-moment assumptions and
the bounds on $H_i,H_0$.

Let $A$ be a traceless Hermitian matrix on a finite input eigenblock.
Then $\rho_i+tA$ is a state for all sufficiently small positive and
negative $t$. Its fourth moment is finite: each selected eigenvector
has finite seventh moment by the construction in
Lemma~\ref{lem:gibbs-regularity}. Finite-dimensional differentiation
on that block gives
\[
 D(\rho_i+tA\Vert\rho_i)
 =\frac{t^2}{2}\mathcal Q_i(A)+o(t^2),\qquad
 \mathcal Q_i(A)=\sum_{l,r}
       \frac{|A_{lr}|^2}{\ell(p_{i,l},p_{i,r})}.
\]
For example, the derivative of the logarithm is the integral of
$(\rho_i+uI)^{-1}A(\rho_i+uI)^{-1}$ over $u>0$ on this block.

To extract a dual lower bound without restricting the actual output
perturbation to a finite eigenblock, we test the Gibbs variational formula
with a centered output observable. Let $Z$ be a Hermitian centered output
test as above. The Gibbs
inequality gives
\[
 D(\widetilde\rho_0\Vert\rho_0)
 \ge t\Tr\widetilde\rho_0Z
    -\log\Tr\exp(\log\rho_0+tZ).
\]
The logarithm of this partition function has first derivative zero
and second derivative $\|Z\|_0^2/h_0$ at zero. This follows from
finite-block exponential differentiation: before centering, only a
finite eigenblock changes, and centering merely multiplies the
exponential by a scalar. Writing $Z_t=\Tr e^{\log\rho_0+tZ}$ and
$\sigma_t=Z_t^{-1}e^{\log\rho_0+tZ}$, the same block decomposition
gives $0<Z_t<\infty$ and makes $\sigma_t$ faithful.
Since $-\log\sigma_t=H_0-tZ+\log Z_t$, the finite energy of
$\widetilde\rho_0$, the bound $H_0\le CW_0$, and boundedness of $Z$
give finite entropy and finite cross entropy against $\sigma_t$.
Lemma~\ref{lem:gibbs-variational} therefore gives the displayed lower bound.
Since the output perturbation induced by
$tA$ is $(1-\kappa)t$ times the raw channel perturbation, the lower
bound has expansion
\[
 t^2\left((1-\kappa)\Tr A T_iZ
             -\frac{\|Z\|_0^2}{2h_0}\right)+o(t^2).
\]
Put $\alpha=(1+\varepsilon)w_i h_0/h_i$. Comparing with
\eqref{eq:relative-entropy-difference} yields
\begin{equation}\label{eq:hessian-duality}
 2(1-\kappa)\Tr A T_iZ-\alpha\mathcal Q_i(A)
 \le\frac{\|Z\|_0^2}{h_0}.
\end{equation}

We can optimize the left side over the completed traceless Hermitian
space for $\mathcal Q_i$. Indeed, trace is continuous in this norm,
since
\[
 |\Tr A|^2\le\left(\sum_l p_{i,l}\right)
                   \sum_l\frac{|A_{ll}|^2}{p_{i,l}}
 \le\mathcal Q_i(A).
\]
Finite eigenblock truncation followed by correcting the trace at one
fixed diagonal entry proves density in the trace-zero subspace.
The constrained dual of $\mathcal Q_i$ is the logarithmic-mean norm
of the observable after subtracting its $\rho_i$ expectation. Thus
the supremum in \eqref{eq:hessian-duality} equals
\[
 \frac{(1-\kappa)^2}{\alpha h_i}
   \|T_iZ-\langle T_iZ\rangle_iI\|_i^2.
\]
It follows that
\[
 (1-\kappa)^2
  \|T_iZ-\langle T_iZ\rangle_iI\|_i^2
 \le(1+\varepsilon)w_i\|Z\|_0^2.
\]
Dividing by $\lambda_i$ gives \eqref{eq:hessian}. The estimate
extends to complex tests by writing real and imaginary Hermitian
parts: symmetry of the logarithmic-mean weights and preservation
of adjoints make each squared norm the sum of the corresponding
two squared norms.

Density now gives the continuous extension. To identify it on a
bounded $Z$, approximate by $P_mZP_m$, where $P_m$ are increasing
finite spectral projections of $\rho_0$, and center each approximant.
They converge in both output arithmetic norms and hence in the
logarithmic-mean norm. The slice is completely positive and unital,
so its Schwarz inequality bounds the corresponding input arithmetic
norms by the raw output norms. Those in turn are at most
$(1-\kappa)^{-1}$ times the modified output norms, because
$\rho_0\ge(1-\kappa)\gamma$. The sliced approximants therefore
converge to the bounded slice in \eqref{eq:slice-map}, as required.
\end{proof}
